\subsection{Definition}

PROPOSITION 1. --- Let A be a commutative Banach algebra. The following conditions are equivalent:

\begin{enumerate}
\def\labelenumi{(\roman{enumi})}
\item
  The weak topology and the Jacobson topology on \(\mathsf{X}(\mathsf{A})\) coincide;
\item
  For every \(\chi \in \mathsf{X}(\mathrm{A})\) and every weakly closed subset \(\mathrm{F}\) of \(\mathsf{X}(\mathrm{A})\) such that \(\chi \notin \mathrm{F}\) , there exists an \(x \in \mathrm{A}\) such that \(\mathcal{G}(x)\) is equal to 1 at \(\chi\) and to 0 on \(\mathrm{F}\) ;
\item
  For every weakly compact subset K and every weakly closed subset F of X(A) such that K ∩ F = ∅, there exists an element \(x \in A\) such that \(\mathcal{G}(x)\) is equal to 1 on K and to 0 on F.
\end{enumerate}

Let \(M \subset \mathsf{X}(A)\) . Saying that M is closed for the Jacobson topology means that, for every \(\chi \in \mathsf{X}(A) - \mathsf{M}\) , there exists an \(x \in A\) such that \(\mathcal{G}(x)\) vanishes on M but not at \(\chi\) (lemma 2 of I, p.~39). Condition (ii) therefore means that every subset of \(\mathsf{X}(A)\) weakly closed is closed for the Jacobson topology, which shows that (ii) \(\Longrightarrow\) (i). Moreover, (iii) \(\Longrightarrow\) (ii) since the set \(\{\chi\}\) is weakly compact in \(\mathsf{X}(A)\). Finally, (i) \(\Longrightarrow\) (iii) by the corollary of prop. 15 of I, p.~81.

DEFINITION 1. --- Let A be a commutative Banach algebra. It is said to be regular if it satisfies the equivalent conditions of Proposition 1.

Note. --- Let \(\tilde{\mathsf{A}}\) be the Banach algebra obtained from A by adjoining an identity element \(e\) . Condition (ii) of Prop. 1 shows that if \(\tilde{\mathsf{A}}\) is regular, then A is regular. Suppose A is regular and let us show that \(\tilde{\mathsf{A}}\) is regular. Consider subsets F and F′ of \(\mathsf{X}(\tilde{\mathsf{A}})\) which are disjoint and weakly closed (hence weakly compact), and let us construct an \(x \in \tilde{\mathsf{A}}\) such that \(\mathcal{G}(x)\) vanishes on F, and is equal to 1 on F′. Let \(\chi_0 \in \mathsf{X}(\tilde{\mathsf{A}})\) be the zero character on A. If \(\chi_0 \notin \mathrm{F}'\) , by condition (iii) of prop. 1 and the hypothesis on A, there exists an element \(x \in \mathsf{A}\) such that \(\mathcal{G}(x)\) vanishes on F and equals 1 on F'. If \(\chi_0 \in \mathrm{F}'\) , one has \(\chi_0 \notin \mathrm{F}\) ; there therefore exists an element \(y \in \mathsf{A}\) such that \(\mathcal{G}(y)\) vanishes on F' and equals 1 on F. The element \(x = e - y\) of \(\tilde{\mathsf{A}}\) then has the required property.

Examples. --- Let us return to the examples of no. 2 of I, p.~17.

The algebra of continuous complex-valued functions tending to 0 at infinity on a locally compact space X (example 3 of I, p.~17) is regular (cf.~I, p.~36, example 1).

The algebra of functions \(n\) times differentiable on \([0, 1]\) (example 4 of I, p.~18) is regular (cf.~I, p.~36, example 2).

If G is a commutative locally compact group and \(\mu\) be a Haar measure on G, then the algebra \(L^{1}(G,\mu)\) (example 7 of I, p.~19) is regular (cf.~II, p.~219, cor. 2).

The algebra of functions that are continuous in the disk \(|z| \leqslant 1\) and analytic in the interior (example 9 of I, p.~20) is not regular (cf.~I, p.~193, exerc. 6).

PROPOSITION 2. --- Let A be a commutative unital regular Banach algebra. Let \(n \geqslant 1\) be an integer and \((\mathrm{U}_{1}, \ldots, \mathrm{U}_{n})\) be an open covering of \(\mathsf{X}(\mathsf{A})\) . There exist elements \(x_{1}, \ldots, x_{n}\) of A with sum 1 such that \(\operatorname{Supp}(\mathcal{G}(x_{i})) \subset \mathrm{U}_{i}\) for \(i = 1, \ldots, n\) .

We prove the proposition by induction on n. The assertion is valid for n = 1. Suppose that \(n \geqslant 2\) and that the assertion is established for n - 1.

There exists an open covering \((\mathrm{V}_{1},\ldots,\mathrm{V}_{n})\) of \(\mathsf{X}(\mathsf{A})\) such that \(\overline{V}_{i}\subset U_{i}\) for every i. By the induction hypothesis, there exist elements \(x,x_{3},\ldots,x_{n}\in A\) such that \(x+x_{3}+\cdots+x_{n}=1\) and \(\operatorname{Supp}(\mathcal{G}(x))\subset\mathrm{V}_{1}\cup\mathrm{V}_{2}\) , \(\operatorname{Supp}(\mathcal{G}(x_{i}))\subset\mathrm{V}_{i}\) for \(i\geqslant3\) . Let \(\mathrm{K}=\operatorname{Supp}(\mathcal{G}(x))\subset\mathrm{V}_{1}\cup\mathrm{V}_{2}\) . Let \(K_{1}\) (resp. \(K_{2}\) ) be the set of elements of K that do not belong to \(V_{1}\) (resp. \(V_{2}\) ). Then \(K_{1}\) and \(K_{2}\) are disjoint compact subsets of K. Since the Banach algebra A is regular, there therefore exists \(y\in A\) such that \(\mathcal{G}(y)=1\) on \(K_{1}\) and \(\mathcal{G}(y)=0\) on \(K_{2}\) . Then \(\mathcal{G}(xy)\) is zero on \(\mathsf{X}(\mathsf{A})-\mathsf{K}\) and on \(K_{2}\) , hence \(\operatorname{Supp}\mathcal{G}(xy)\subset\overline{\mathrm{V}}_{2}\subset\mathrm{U}_{2}\) . Likewise, \(\mathcal{G}(x(1-y))\) is zero on \(\mathsf{X}(\mathsf{A})-\mathsf{K}\) and on \(K_{1}\) , hence \(\operatorname{Supp}\mathcal{G}(x(1-y))\subset\overline{\mathrm{V}}_{1}\subset\mathrm{U}_{1}\) . The elements \(x_{1}=x(1-y)\) , \(x_{2}=xy\) , and \(x_{3},\ldots,x_{n}\) then satisfy the properties of the proposition.

COROLLARY 1. --- Let A be a commutative unital regular Banach algebra, let I be an ideal of A, and let \(f: \mathsf{X}(\mathsf{A}) \to \mathbf{C}\) be a continuous function. Suppose that, for every \(\chi \in \mathsf{X}(\mathsf{A})\) , there exists an element \(y_{\chi} \in \mathrm{I}\) such that \(f = \mathcal{G}(y_{\chi})\) in a neighborhood of \(\chi\) . Then there exists an element \(y \in \mathrm{I}\) such that \(f = \mathcal{G}(y)\) .

Since X(A) is compact, there exists a finite open covering \((\mathrm{U}_{1},\ldots,\mathrm{U}_{n})\) of X(A), and elements \(y_{1},\ldots,y_{n}\) of I such that \(f=\mathcal{G}(y_{i})\) on \(U_{i}\) . By prop. 2, there exist elements \(x_{1},\ldots,x_{n}\) of A with sum 1 such that \(\operatorname{Supp}(\mathcal{G}(x_{i}))\subset\mathrm{U}_{i}\) for every i. Let \(y=x_{1}y_{1}+\cdots+x_{n}y_{n}\) . This is an element of I that has the required property. Indeed, let \(\chi\in\mathrm{X}(A)\) . For \(1\leqslant i\leqslant n\) , we have \(\mathcal{G}(x_{i})(\chi)\mathcal{G}(y_{i})(\chi)=\mathcal{G}(x_{i})(\chi)f(\chi)\) since \(\mathcal{G}(y_{i})(\chi)=f(\chi)\) if \(\chi\in U_{i}\) , and \(\mathcal{G}(x_{i})(\chi)=0\) if \(\chi\notin U_{i}\) . It follows that

\[
\mathcal {G} (y) (\chi) = \sum_ {i = 1} ^ {n} \mathcal {G} (x _ {i}) (\chi) \mathcal {G} (y _ {i}) (\chi) = f (\chi) \sum_ {i = 1} ^ {n} \mathcal {G} (x _ {i}) (\chi) = f (\chi).
\]

COROLLARY 2. --- Let A be a commutative regular Banach algebra, I an ideal of A, and \(f\colon\mathsf{X}'(\mathrm{A})\to\mathbf{C}\) be a continuous function. Suppose that, for every \(\chi\in\mathsf{X}'(\mathrm{A})\) , there exists an element \(y_{\chi}\in\mathrm{I}\) such that \(f=\mathcal{G}'(y_{\chi})\) in a neighborhood of \(\chi\) . Then there exists an element \(y\in\mathrm{I}\) such that \(f=\mathcal{G}'(y)\) .

Let \(\tilde{\mathsf{A}}\) be the Banach algebra obtained from \(\mathsf{A}\) by adjoining an identity element. Then \(\tilde{\mathsf{A}}\) is regular (remark 1), and \(\mathsf{X}'(\mathsf{A}) = \mathsf{X}(\widetilde{\mathsf{A}})\) ; it therefore suffices to apply cor. 1 to \(\tilde{\mathsf{A}}\) and to the ideal I.

If I is an ideal of a commutative Banach algebra, recall (cf.~I, p.~30) that we denote by V(I) the set of \(\chi \in \mathsf{X}(\mathrm{A})\) whose kernel contains I, in other words the set of \(\chi \in \mathsf{X}(\mathrm{A})\) where all the functions \(\mathcal{G}(x)\) for \(x \in \mathrm{I}\) vanish. This is a subset of \(\mathsf{X}(\mathrm{A})\) closed for the Jacobson topology.

PROPOSITION 3. --- Let A be a commutative regular Banach algebra, I an ideal of A, and K a compact subset of X(A) disjoint from V(I). There exists an element \(x \in I\) such that \(\mathcal{G}(x) = 1\) for every \(x\) in K.

This is a special case of prop. 15 of I, p.~81, taking into account the fact that the Jacobson topology coincides with the weak topology on X(A).
% semantic-fragments: legacy-input-seam
\relax{}
